\documentclass[10pt]{article}
\usepackage[a4paper,margin=23mm]{geometry}
\usepackage{amsmath,amssymb,booktabs,graphicx,hyperref}
\setlength{\emergencystretch}{2em}
\title{Query-Conditioned Evidence Fusion for Reliable 3D Scene Responses}
\author{TMM working manuscript --- controlled validation in progress}
\date{October 17, 2026}
\begin{document}
\maketitle
\begin{abstract}
Strong visual teachers provide useful but incomplete evidence for scene queries.
Matching an observed region does not determine whether its surroundings are
background or whether it contains the complete requested object. We study how
query evidence should be admitted and combined into a three-dimensional response.
Our method separates query applicability, spatial evidence authority and
geometric aggregation, with typed semantic and complete-instance readouts. It
preserves native matching without requiring a new common feature representation.
The v10 system is our implementation, evaluated across five outputs in 31 scene
jobs. It obtains 70.40\% LERF semantic IoU, 66.96\% selected-set projection IoU,
61.01\% ScanNet 19-class mesh mIoU and 86.51\% NVOS IoU, with four query modalities
evaluated on UQIS. Historical paired controls connect identity admission,
admitted-response use, local image correspondence and complete membership to
their measured effects. Input and metric differences are disclosed for external
comparisons. Further controls show why lower source loss and response replacement
alone do not ensure better answers. A fixed eight-task NVOS study improves IoU
from 86.51\% to 87.78\% by changing source evidence construction and sampling,
with the original teachers, geometry and output threshold retained. An explicitly
image-conditioned extension that completes and admits reference responses reaches
93.17\% after native-ray 3D writeback, with 98.95\% AP; its direct image output
reaches 93.50\%. This is the fixed source-quality secondary policy; the declared
source-feasibility primary reaches 92.94\% 3D IoU. Both are measured on all eight
tasks, with original negative scribbles and frozen DINOv2 explicitly added.
These are separate access
conditions, not a single input-identical improvement. The framework has effective typed
implementations; one learned rule with gains across every query type is not
established by the present evidence.
\end{abstract}
\section{Introduction}
A scene query asks for an identity and its spatial extent. Text may name an
object, an image may show one of its parts, and a point may identify one visible
surface. A useful three-dimensional response must connect that evidence to the
requested object while rejecting distractors and recovering members beyond the
observed region. Strong teacher features provide valuable correspondences, but
a high match score alone does not specify the extent of a correct answer.

Two decisions are easily conflated. First, an observation may or may not refer
to the requested identity. Second, even a correctly matched observation may
support only part of that identity's extent. Treating its exterior as background
can erase unseen object parts; accepting an unrelated region can propagate
false positives through otherwise accurate geometry. The problem is therefore
not only how to aggregate evidence, but which evidence is entitled to support
or exclude each part of the response.

We study \emph{query-conditioned evidence fusion}: native observations are
admitted according to their applicability to the query, assigned spatial
support and exclusion authority, and combined through a fixed scene carrier.
Rejected or unavailable evidence remains unknown. Text applicability uses
candidate-relative native matching; reference tracks, categorical competition
and explicit points retain their respective conditioning semantics. This
response-level interface uses existing teacher representations without learning
a common latent feature field.

Evidence formation must also preserve the intended output capability.
Continuous semantic responses and complete selected instances have different
optimization objectives and readouts. We retain both, using geometric response
fusion for semantic evidence and complete-membership mechanisms for actual
selection. This avoids interpreting silhouette membership as a calibrated
semantic probability, or assuming that a better score field automatically
improves an independent selected set. The shared evidence contract has typed
implementations; it does not require an identical solver for every query type.

Our implementation, recorded as v10 for reproducibility, covers LERF semantic
segmentation and selected-set projection, ScanNet open-vocabulary segmentation,
NVOS reference segmentation and multimodal UQIS. Controlled component studies
connect identity admission, admitted-response use, local image correspondence
and complete membership to measured improvements. Further matched experiments
separate changes in evidence construction, reconstruction, response decisions
and output interpretation. These controls also expose cases where better source
fit fails to improve target answers. External comparisons retain their declared
teacher, geometry and input-access differences.

Our contributions are threefold. First, we formulate a query-conditioned
evidence contract that separates identity applicability, spatial support,
exclusion and missing information. Second, we implement this contract in a
scene-query method that combines native evidence through geometry while
preserving both semantic response and complete selection. Third, we evaluate
five task outputs with matched component controls, protocol checks and complete
positive and negative feedback. Optional response-replacement and NVOS
extensions are evaluated separately from the established system, so their
measured effects can be attributed to the operations they actually change.

\section{Related work}
\paragraph{Language and visual evidence in Gaussian scenes.}
VALA studies visibility-aware semantic lifting and robust multi-view
aggregation~\cite{vala}. Occam's LGS~\cite{occam} provides a Gaussian language
baseline. LUDVIG lifts features or masks and uses graph diffusion for scene
queries~\cite{ludvig}. Their published results and our retained local
reproductions provide comparison context. We do not claim visibility weighting,
robust aggregation or native matching individually as new operations.

EviSplat~\cite{evisplat} preserves observation features within instances and
combines query-selected evidence with local Gaussian support. Query-dependent
evidence use is therefore established. Our focus is the authority of a
query-conditioned observation to support or exclude a response, and its
interaction with complete extent across typed tasks. Matched controls establish
measured contributions. Teacher and geometry differences limit causal
attribution of gains over published methods.

\paragraph{Response verification.}
Our optional complete-response selector predicts actual answer utility from
observable evidence and uses a shared acceptance margin. Consistent ordering
and conservative retention do not guarantee accurate transfer. This extension
is evaluated with fixed candidates and query/retention controls; it is distinct
from the fusion and membership that produce our v10 answers.

\section{Query-conditioned evidence fusion}
\label{sec:method}
The method forms a response in three stages: determine which observations
apply to the query, assign their spatial support and exclusion authority, and
combine their admitted responses through geometry. Semantic and complete-instance
outputs retain their own objectives and readouts. The following equations
describe the implemented system and identify its query-specific realizations.

\begin{figure}[ht]
\centering\includegraphics[width=\linewidth]{figures/method_overview.pdf}
\caption{Our evidence-to-response method. Applicability and spatial authority
determine how native evidence acts on the response. Geometric aggregation and
complete membership have distinct typed implementations. Actual selections are
evaluated separately from continuous semantic responses.}
\label{fig:method}
\end{figure}

\subsection{Native evidence and query applicability}
For query $q$, observations $\mathcal E_q$ supply responses $y_e\in[0,1]$,
nonnegative geometric operators $A_e$ and visible domains $v_e$. Nonempty
operator rows are normalized. Text similarities and local visual correspondence
stay in their native teacher spaces; points and reference masks retain their own
semantics. We combine response evidence, not raw feature vectors. LERF, NVOS
and ScanNet use declared source-assisted observations; UQIS reads its map without
source RGB or depth at query time.

Let $g_e(q)$ express whether observation $e$ applies to the requested target.
For text candidates, we compare one- and two-component native-score models
within candidates already detected for the same query. If their Bayesian
information criteria are $B_1,B_2$, the implemented support is
\begin{equation}
 a_e=(1-\pi_2)+\pi_2\gamma_{\rm high}(s_e),\quad
 \pi_2=\sigma((B_1-B_2)/2),\quad g_e=\mathbf1[a_e\geq0.5],
 \label{eq:admission}
\end{equation}
where $\gamma_{\rm high}$ is responsibility for the higher-score component.
A single mode supplies no relative identity veto. Degenerate or missing evidence
does not invent one. Rejected regions become unknown rather than background.
Mixture support is not a calibrated probability of semantic truth.

This rule has a controlled LERF result. Reference queries instead inherit
identity through a positive-prompt mask and source tracks; categories combine
native probabilities with literal-category evidence. Image queries use local
correspondence to select identity hypotheses. Actual points supply explicit
spatial conditions. These are typed estimators, not unreported copies of
Equation~\ref{eq:admission}.

\subsection{Spatial authority and geometric response formation}
An applicable observation need not describe complete extent. Let $p_e,n_e$
denote disjoint support and exclusion domains and $c_e^+,c_e^-$ their weights:
\begin{equation}
 w_e=v_e g_e(q)(c_e^+p_e+c_e^-n_e).
 \label{eq:authority}
\end{equation}
Unobserved space and rejected detections carry no exclusion solely because they
lack a match. Predicted reference-mask complements and category competition can supply
exclusion under their stated completeness assumptions. These are fallible
provider estimates, not hard negative user prompts. Our LERF path vetoes negatives where
native positive support exists. NVOS uses interior positive/negative regions and
tracked negative reliability. ScanNet keeps both native priors and source weights.

For direct consensus, define $m=\sum_e A_e^\top w_e$ and
$b=\sum_e A_e^\top(w_ey_e)$. The response is
\begin{equation}
 r_i=\begin{cases}
 (b_i+\lambda_i\bar r_i)/(m_i+\lambda_i),&m_i>10^{-8},\\
 r_{0,i},&\text{otherwise}.
 \end{cases}
 \label{eq:consensus}
\end{equation}
LERF uses $\lambda=0$. ScanNet uses a two-native mean $\bar r$ and prior strength
0.1 times visible mass. This is an established weighted-consensus operation;
the tested design concerns admitted evidence and its actual use in the response.
The active ScanNet implementation has no fitted local residual head.

NVOS instead minimizes robust rendered discrepancy,
\begin{equation}
 \min_{r\in[0,1]^N}\sum_e\langle w_e,\rho_{0.25}(A_er-y_e)\rangle,
 \label{eq:reference}
\end{equation}
with 120 mass-preconditioned updates from a neutral 0.5 response. Its output
adapter zeros unsupported elements. Changing this solver while preserving
evidence is an ablation. Our controls do not support replacing every typed
operator by one inverse solver.

\subsection{Complete extent and typed outputs}
For source-assisted actual selections, our implementation also fits membership
against the silhouette of the selected Gaussians. This is distinct from rendering
their semantic scores in the full scene. Let $z_i\in[0,1]$ be membership and
$o_{vi}(u)$ the uncomposited per-Gaussian opacity at pixel $u$. The selected
silhouette has the compositing form
\begin{equation}
 \alpha_v(u;z)=1-\prod_i(1-z_i o_{vi}(u)).
\end{equation}
The implemented loss combines foreground/background-balanced binary
cross-entropy and soft Dice on admitted source pixels, plus a 0.1-weighted
projected-footprint membership penalty. The footprint is normalized by its
cost at source-initialized membership. Optimization uses at least 240 Adam
updates (extended to the source-view count when larger) at learning rate 0.05,
cycling through the active views. Actual membership is then
thresholded at 0.5 and its selected opacity is rendered. The existing projection
readout thresholds quantized 8-bit alpha above 10. This describes an existing
v10 capability, not a new solver introduced by the subsequent NVOS experiment.

\paragraph{Native image identity and conditional complete membership.}
For an image query, let $q_u,f_v$ be normalized query and regional patch
features, with unit-sum weights $a_u,b_v$. Local correspondence uses the
symmetric weighted Chamfer score
\begin{equation}
 k(Q,R)=\tfrac12\left[\sum_{u\in Q}a_u\max_{v\in R}q_u^\top f_v+
                         \sum_{v\in R}b_v\max_{u\in Q}q_u^\top f_v\right].
\end{equation}
The deployed image route averages the candidate ranks from native SigLIP
matching and local DINOv2 correspondence. It selects a nonempty region, then
uses its actual member nearest the coordinate-wise 3D median as the identity
anchor. Both teacher spaces remain unchanged. An actual point supplies a
geometric anchor directly; an image-derived anchor remains an identity estimate.

For anchor $a$ and every carrier element $i$, the complete-membership decoder
predicts
\begin{equation}
 p_i=f_\theta(\phi(i,a)),\qquad S_a=\{i:p_i>\tau\},\quad p_a=1.
\end{equation}
The 15-dimensional evidence vector retains view-balanced nested-mask votes,
any-mask/all-mask agreement and their disagreement, common visibility and its
presence flag, anchor/element visibility fractions, bounded signed displacement
and distances, native appearance agreement, and native-support availability.
Nested masks share a view's vote instead of becoming independent observations.
The decision is evaluated over the complete carrier, without requiring an
intermediate mask to be a pure or complete instance.

The deployed decoder is an Extra Trees classifier with 64 trees, at most 512
leaves per tree and minimum leaf size 32. It is inherited from training on
independent 3D instance labels in 22 real-observation rooms; three separate
calibration rooms select $\tau=0.35$. Neither split overlaps the benchmark
rooms. This supervised decoder is disclosed as part of the method; freezing
teacher features does not make the complete system training-free. The same
anchor-positive rule is applied to inferred image anchors. Its identity
condition is fallible, and an ablation removing only that assignment leaves
all measured selected sets unchanged.

Continuous semantic response, category argmax and actual instance selection
have different readouts. LERF and NVOS maintain independent Gaussian sets;
changing the semantic response does not change that set automatically.
Selected-set projection IoU uses 2D GT, not volumetric GT. UQIS complete
membership is not Equation~\ref{eq:consensus}. These distinctions are essential
to evaluating functional improvement.

\paragraph{Why the optimization target and readout must agree.}
Removing unselected Gaussians changes ray transmittance, not just the values
averaged by the original renderer. A simple ray shows the distinction: an
unselected front Gaussian with opacity 0.9 followed by a selected Gaussian with
opacity 0.2 gives a full-scene normalized member response
$0.02/(0.9+0.02)\simeq0.0217$. Rendering the selected set alone instead gives
alpha 0.2. The two outputs disagree under their existing thresholds despite
identical membership. This is an illustrative counterexample to treating member
scores as calibrated semantic responses, not an accuracy result. Accordingly,
functional inheritance preserves evidence, optimization target and intended
readout together. Our field-by-readout experiment tests this distinction on
retained NVOS fields.

\paragraph{Image-conditioned response refinement.}
For an observed view $t$, let $z_t^0=A_t s^0$ be the rendered v10 response.
We use its thresholded extent as a box and its bounded log-odds as a dense
prompt for a frozen image decoder $D$:
\begin{equation}
 p_t=D(I_t,\operatorname{box}[z_t^0\geq0.5],
       \operatorname{resize}[\operatorname{logit}(\operatorname{clip}(z_t^0))]).
\end{equation}
Clipping uses $[10^{-4},1-10^{-4}]$; the dense prompt is resized to the
decoder's actual $288\times288$ grid. The box and dense mask are fallible
predictions, not extra annotations. Single-mask decoding avoids target-label
candidate selection. Paired reconstructed/observed RGB controls isolate image
access under this fixed interface. Sparse-point and unboxed controls test
whether the inherited response preserves identity while restricting extent.

The coarse control uses confidence-weighted area pooling of $p_t$ on a
$512\times384$ grid. Native-ray writeback instead retains $y_t=p_t$ on the
$1008\times756$ output grid and uses $w_t=|2p_t-1|v_t$. Here $v_t$ admits rays
with nonzero observed mass and more than half of their rendered alpha supplied
by the original observed Gaussians. The normalized observation matrix
$\bar A_t$ is formed on the same ray grid, using the fixed original Gaussian
geometry. Initialize at $s^0$ and take 120 bounded Huber updates for
\begin{equation}
 \min_{0\leq s\leq1}\sum_u w_{tu}\rho_{0.25}((\bar A_t s)_u-y_{tu}),
 \qquad s_i=s_i^0\ \text{when no admitted observation supports }i.
\end{equation}
We measure both $p_t$ and the actual re-rendered 3D response. Geometry and
teacher weights remain fixed. Because this update uses view $t$'s image,
the result is view-conditioned; a query-global field and preserved source-view
accuracy do not follow from target-view gains. Original source-teacher
residuals are reported as a separate compatibility check.

\paragraph{Completing query scope before transport.}
The current NVOS extension builds a complete reference response before cross-view
transport. For reference image $I_r$, let $\mathcal P,\mathcal N$ be original
positive and negative scribble patches. A patch containing both signs is unknown;
unmarked patches do not supply negative labels. With frozen normalized DINOv2
descriptors $f_r$, define
\[
 h_r(u)=\max_{v\in\mathcal P} f_r(u)^\top f_r(v)
       -\max_{v\in\mathcal N} f_r(u)^\top f_r(v),
 \qquad a_r(u)=\sigma(h_r(u)/\tau_r).
\]
The source-only scale $\tau_r$ is the median signed leave-one-out reference
margin divided by $\log9$. Missing signed support or a nonpositive scale falls
back to the original reference mask. Upsample $a_r$ and use the same dense-mask
and extent-box interface to obtain a semantic reference response from frozen
SAM3. A second candidate additionally uses 16 deterministically spread points
per original scribble sign. Threshold both at 0.5; for the constrained candidate,
impose all original positive/negative pixels exactly, without relabelling
unknown pixels. Agreement with sparse annotations does not prove complete scope.

Two fixed reference-only admission policies are compared. The feasibility policy
retains the semantic candidate when it satisfies every original signed pixel,
otherwise choosing the constrained candidate. The source-quality policy chooses
\[
 \hat k=\arg\max_{k\in\{\mathrm{semantic},\mathrm{constrained}\}}
             \widehat Q_r(k),
\]
where $\widehat Q_r$ is SAM3's predicted mask quality at reference decoding,
before the final hard annotation constraints. Ties prefer the semantic candidate;
missing quality falls back to feasibility. These scores are fallible predictions,
not target GT or calibrated correctness probabilities. Neither rule reads a
scene identifier, target score or task-specific threshold.

The selected whole reference mask initializes frozen two-frame SAM3 memory,
producing $y_t=T(I_r,m_r^{\hat k},I_t)$. The same native-ray fit above then
creates and actually re-renders the 3D response. Original negative scribbles,
DINOv2 and observed target RGB remain declared extensions to v10's input contract.
This changes response construction and admission, keeping teacher weights and
Gaussian geometry fixed. Both fixed policies are reported on all eight tasks;
reference-admission design follows the opened-cohort completion diagnosis.
Earlier geometric region-admission and paired-residual alternatives remain
matched diagnostics, rather than additional required stages of this pipeline.

\paragraph{Optional complete-response replacement.}
A subsequent extension generates intact alternatives and predicts relative task
utility from query/response signatures. A calibrated margin retains the
incumbent or returns one complete alternative. Its L11 variant uses a structural
potential plus a query-conditioned residual. This extension is evaluated
separately: its small image gain and unfinished LERF/NVOS transfer do not account
for the historical component gains.

\section{Experiments}
\label{sec:experiments}
\subsection{Datasets, metrics and comparison protocol}
We use one fixed in-house method recipe across 31 scene jobs. Four LERF scenes yield
two outputs, so the cohort comprises five evaluation tasks rather than five
independent datasets. All reported numbers are percentages. The existing
benchmark rooms have been used during development; they are not an untouched
final test. Scene-level intervals describe this cohort rather than certify
generalization after repeated research feedback.

\begin{table}[ht]
\centering\small
\begin{tabular}{lp{.25\linewidth}p{.43\linewidth}}
\toprule
Task & Cohort and primary metric & Input and evaluation contract \\
\midrule
LERF-2D & 4 scenes, 208 observations; mIoU and localization & Source-assisted text; semantic response rendered before thresholding. \\
LERF-3D selection & Same scenes; projection mIoU and Acc.@0.25 & Independent selected Gaussians rendered to 2D GT; not volumetric GT IoU. \\
ScanNet OVS & Paper8 Gaussian mIoU/mAcc; internal10 mesh mIoU & Fixed 19/15/10 menus; Gaussian and mesh domains reported separately. \\
NVOS & 8 tasks, 7 scene clusters; IoU & Retained positive-reference, cached-source variant; external prompt/RGB access declared. \\
UQIS & 9 rooms, 67 instances, 268 queries; set IoU and AP & Text, image, 2D point and 3D point; map-only query-time access. \\
\bottomrule
\end{tabular}
\caption{Evaluation units. LERF and UQIS average within scenes then across
scenes; NVOS averages eight tasks, with three seeds per task for LUDVIG.}
\label{tab:protocols}
\end{table}

\paragraph{Published results and verified reproductions.}
We reuse values reported in VALA Tables 1--2 and LUDVIG Tables 1/D, including
their cited baseline results~\cite{vala,ludvig}. No new external baseline training
is needed. Every local reproduction is linked to its frozen result hashes;
four-scene, eight-scene and eight-task aggregates are recomputed from the
original reports. Table origins are \textbf{P}: published, \textbf{R}: locally
reproduced, \textbf{O}: our implemented method (v10 or a named internal control), and \textbf{M}: our method with a named measured extension.
Missing extension results remain pending in the registry; they do not erase the evaluated v10 method results.
We do not bold a global winner across different input contracts.

\subsection{Implementation and supervision}
The fixed v10 recipe includes native teacher matching, source evidence, fixed
geometry and typed task readouts. SAM3~\cite{sam3} source masks and native
SigLIP2~\cite{siglip2} scores support LERF. ScanNet combines original SigLIP2
and Mask-Adapter~\cite{maskadapter} priors with
literal-category SAM3 evidence. NVOS uses positive-prompt SAM3 reference tracks
and a robust rendered response. All nine UQIS maps use SigLIP2 giant and
DINOv2 ViT-B/14~\cite{dinov2} for image matching;
complete membership inherits supervision on 22 labelled rooms and calibration
on three other rooms. These room identities are disjoint from benchmark rooms.
The original feature fields and teacher weights are not retrained in this study.

The optional L11 selector additionally uses independent text/image/category
utility supervision. Its 22 fit, three calibration and four development rooms
provide 629/207/175 text expressions, 237/28/25 image queries and 116 category
episodes. Only 88 fit-room category episodes enter training; calibration and
development use full maps. Query-image views and neighbors are excluded from
their candidate observations; full-scan geometry remains available. One margin
(0.15052) protects 15 calibration room/output cells. Five room folds fit a
64-tree structural prior plus a 64-tree query residual. All 19 ScanNet/UQIS jobs
are predicted before scoring. These are extension settings, not prerequisites
for the v10 main results. Development and calibration informed model design.

\subsection{Comparison with existing methods}
\paragraph{LERF semantic segmentation and projected selection.}
Tables~\ref{tab:lerf2d}--\ref{tab:lerf3d} reuse the published context and verified
Occam/VALA reproductions. The VALA reproduction follows its semantic and
evaluation path on compatible Occam RGB geometry, rather than retraining the
entire original geometry pipeline. The v10 values are our method results. The optional shared replacement has
not yet been measured on these two outputs. In particular, semantic-response improvements cannot be substituted
for independent selected-set improvements.
Our method includes SAM3 text-conditioned source evidence and native SigLIP2
matching; external VALA scores are not matched-teacher controls. An independent
CPU audit exactly reproduces all 208 retained selected-set mask IoUs. Semantic
IoU/localization records are reaggregated, without a fresh semantic render in
that audit.

\begin{table}[ht]
\centering\small\input{tables/comparison_lerf2d.tex}
\caption{LERF-2D. P values are reported in VALA Table 1~\cite{vala}.
Localization is peak-in-box accuracy, not ranking AP. O uses source-assisted
full teachers; these are v10 method results, not new replacement gains.}
\label{tab:lerf2d}
\end{table}
\begin{table}[ht]
\centering\small\input{tables/comparison_lerf3d.tex}
\caption{LERF selected-Gaussian projection. P values are reported in VALA
Table 1~\cite{vala}; Acc.@0.25 counts projected masks with IoU above 0.25.
VALA's published value is 58.02/82.85; 43.29/64.30 belongs to Dr. Splat.
The local VALA R row uses compatible RGB geometry.}
\label{tab:lerf3d}
\end{table}

\paragraph{ScanNet open-vocabulary segmentation.}
Table~\ref{tab:scannet} separates the paper-eight Gaussian metric from the
original ten-room mesh metric. For the metric bridge, we score unchanged v10
labels with the existing VALA evaluator: anisotropic Mahalanobis pseudo labels,
five-scale radius, top-1000 class-balanced votes and opacity--volume weights.
All eight scene predictions are hashed before opening their target labels.
This aligns the metric and scene list, but not the input carrier: v10 uses
provided unlabelled mesh geometry, while VALA uses optimized RGB Gaussians.
The bridge is a protocol diagnostic, not an additional learned gain. The
replacement policy retains every original menu on the internal ten-room track.
Independent scoring reproduces all ten mesh scenes and all eight Gaussian
scenes. All 1,134,207 pseudo labels agree, including 256 checks with the released
VALA code. On the same eight rooms, 19-class mesh mIoU is 61.0567 versus
66.9480 for Gaussian scoring: this 5.8913-point change is metric-dependent.

\begin{table}[ht]
\centering\small
\resizebox{\linewidth}{!}{\input{tables/comparison_scannet.tex}}
\caption{ScanNet OVS. P rows from VALA Table 2~\cite{vala}, R and the metric
bridge use paper8 Gaussian scoring. Rows marked 10-mesh use the original
ten-room mesh evaluation and cannot be ranked against the paper8 rows.
The metric bridge still differs in geometry and teacher access.}
\label{tab:scannet}
\end{table}

\paragraph{NVOS.}
Table~\ref{tab:nvos} includes the complete three-seed LUDVIG-SAM reproduction.
Its positive scribbles, all-view RGB geometry and online target-RGB access
follow the released implementation. Our original v10 method uses a
positive-derived reference SAM mask and cached source tracks, without online
target RGB. These are contextual method comparisons with declared access
differences. Later image-conditioned extensions explicitly permit observed
target RGB; they retain the original geometry and are reported separately.
The audit also finds a GT boundary convention difference: our retained scorer
uses $g>127$, while LUDVIG uses integer division by the mask maximum. Applying
only its GT convention changes our score from 86.5134 to 86.6835; historical
values remain unchanged. Switching bilinear implementations changes less than
0.00002 points. Borrowing its min--max and fixed 0.25 response threshold yields
86.6085 on our GT convention, and is not a full LUDVIG reproduction. Flower,
trex and orchids account for 78.04\% of the gap to the local 91.2577 reproduction.
Flower/orchids primarily miss foreground; trex also activates background.
N04 pairs reconstructed and observed RGB at the same camera, with identical
response-derived prompts and frozen image decoder. Observed RGB increases
IoU from 87.5575 to 89.0636 for direct image responses, and from 86.9242 to
88.3201 after actual 3D writeback. This isolates an image-access intervention
within this decoder; it does not explain the entire difference to LUDVIG.
The 0.7435-point loss between direct image and 3D outputs leaves transport,
observation pooling, geometry and the fitting objective to be separated.
The 3D candidate raises IoU by 1.8067 points over v10 but lowers AP, so it is
not a uniform improvement across metrics. Source-teacher residuals also rise;
this is a view-conditioned update, not a demonstrated query-global field.
N05 finds that discarding negative dense hints loses 2.6331 IoU points against
the matched signed-hint control. Sparse hints recover flower but admit another
object in horns-center. Complete prompt controls and both GT conventions are
reported in the diagnostic supplement; no per-task winners are combined.
N07's reference-memory transfer does not exceed the static image-conditioned
mean. N08 fixes all three image-evidence arrays and fits responses on the native
1008-by-756 ray grid instead of 512-by-384. The predeclared sequence-memory
primary arm reaches 88.1124 IoU. The static secondary arm reaches 88.8424,
up 0.5223 points over its coarse control with all eight tasks improved, leaving
a 0.2212-point gap to its unchanged direct image evidence. An accompanying
CUDA arithmetic correction changes matched coarse-grid mean IoU by less than
0.000006 points. These results support resolution-consistent transport, not
a benefit from memory itself. N09 replays this static control with identical
per-task binary counts in both GT tracks and measures 96.5817 AP, below v10's
98.0222. The static image arrays match exactly; the native continuous arrays
are not bitwise identical to N08. Neither run certifies SOTA.

N09 adds frozen DINOv2 descriptors and the official negative reference scribbles
to test complete query transfer beyond the initial mask. Signed-scribble authority
and reference-mask/complement authority use matched inputs and downstream
operations. Their native IoUs are 81.5803 and 86.0812. Signed evidence recovers
missing regions in flower and orchids, but confuses other instances in
horns-center and leaves. Its source leave-one-out accuracy of 99.80--100\%
therefore does not establish target-instance correctness. N10 writes the exact
affinity probabilities directly to 3D, yielding 77.7307 and 84.2098 IoU;
omitting SAM does not recover the static mean. All 72 N09 and 48 N10 predictions
are sealed before their respective scoring. Detailed paired controls, AP and
full-cohort error maps appear in the diagnostic supplement. The added teacher
and negative-prompt access are disclosed; no per-task branch selection is used.

N11 fixes those image proposals and compares unrestricted additive completion,
v10-response anchoring and signed-scribble geometry anchoring. Their actual 3D
IoUs are 83.5805, 89.9673 and 89.9734. The predeclared geometry primary gains
1.1310 points over static native writeback (six tasks improve, two regress),
but loses 0.1540 AP points. Flower improves from 85.9782 to 94.9302, while
instance constraints prevent the unrestricted branch's severe horns-center
and leaves regressions. Orchids remains 73.5362: the rule rejects most of the
proposal's missing-flower recovery. The 0.0061-point difference between the
two anchor policies does not establish superiority of geometric witnesses.
All 72 predictions precede scoring. N11 is a development improvement, not a
SOTA certification or a common improvement across all five benchmarks.
N12 moves the intervention to the reference image. Original, semantic-completed
and signed-constrained reference masks enter identical two-frame memory and
native fitting, giving 88.2002, 88.5226 and 91.4774 IoU. Constrained completion
raises orchids to 94.5224 but harms horns-left by modifying a reference mask
that already agrees with all original scribbles. N13 instead transfers the
paired memory difference around the static response: probability and log-odds
variants give 90.7469 and 89.4416. The probability residual harms all eight tasks
versus its matched complete replacement, so it is not retained as a required
method stage.

N14 fixes two source-only admission rules after this N12 development diagnosis.
Feasibility admission (primary) reaches 92.9424 IoU / 98.8599 AP; source-quality
admission (secondary) reaches 93.1730 / 98.9451. Their LUDVIG-GT IoUs are 92.9852
and 93.2104. Quality exceeds static native writeback by 4.3305 points and N11
by 3.1996, improving five and harming three tasks in either comparison. It
improves only trex against feasibility, while harming fern and horns-center;
five tasks tie. Thus the aggregate does not establish a universally better
quality criterion. Flower, orchids and horns-left reach 96.7139, 95.6011 and
94.8079. The quality image output is 93.4971, with 0.3241 mean loss after actual
3D fitting; internal background leakage and fine structures remain limiting.

All 72/56/56 N12/N13/N14 outputs respectively precede their GT scoring. N14
freshly replays selected teacher responses and verifies exact parent hashes,
then fits and re-scores actual 3D outputs; it does not substitute cached
per-task maxima. Original negative scribbles, DINOv2 and target RGB are shared
with N11 but extend the original v10 input contract. The other 23 jobs remain
unchanged references. This is adaptive development evidence, without a matched
external SOTA certification or all-five default promotion.

N15 freezes the N14 reference-quality decision and transports both original and
completed references to all 318 authorized source views. Matched source-only
fusion reaches 88.0061 and 92.7461 IoU, respectively: completion gains 4.7400
points with no target-image evidence in either field. Equal-group fusion with
the N14 target response reaches 93.0302 IoU / 98.9797 AP, versus the replayed
N14 control's 93.1730 / 98.9451. It reduces false positives but also true
foreground on every task; precision increases while recall falls. Both GT
conventions retain the same mean-IoU ordering. The predeclared joint candidate
therefore fails the higher-IoU, non-decreasing-AP criterion, and this fusion
branch is stopped. All 48 outputs precede scoring, and the three N14/v10
controls reproduce the original binary counts. The source-only result supports
transfer of reference-completion capability; it does not establish that more
source constraints improve the strongest target-conditioned output.

N16 re-decodes the exact complete N14 target response once, using either the
full target image or a response-derived crop with 25 percent context per side.
The same crop transforms RGB and prior; outside probabilities are retained
before native fitting. Full-frame decoding gives 93.0998 IoU / 98.2626 AP.
The local primary reaches 93.2805 / 98.4663, versus the replayed N14 control's
93.1730 / 98.9450. Thus the highest observed native IoU increases by 0.1075
points, while AP loses 0.4787. Four tasks improve and four regress. Local
decoding exceeds its matched full-frame control by 0.1806 IoU and 0.2037 AP
points, but does not retain the N14 ranking quality. Fern and horns-center
receive full-frame windows under the fixed rule and have identical predictions
in the two controls. Both GT conventions show the tradeoff; omitting the whole
horns scene removes the mean-IoU gain over N14. All 56 predictions precede
scoring. The branch fails its predeclared joint-metric criterion and is stopped;
N16 remains a measured IoU/AP tradeoff, while N14 remains the retained balanced
recipe. Neither result certifies an external same-protocol SOTA ranking.
An exact eight-task score audit (A24, diagnostics) preserves all hard masks and
changes mean image AP by less than 0.002 points under higher-precision probability
interpolation; it does not overturn the N16 stopping decision or alter 3D scores.

\begin{table}[ht]
\centering\small\resizebox{\linewidth}{!}{\input{tables/comparison_nvos.tex}}
\caption{NVOS IoU. Our rows use all eight tasks. P denotes published values from
LUDVIG Tables 1/D~\cite{ludvig}, VCAR Table 1~\cite{vcar}, LEGO Table 1~\cite{lego}, WildSeg3D Table 1~\cite{wildseg3d}, and Trace3D Table 4~\cite{trace3d}.
VCAR uses seven targets; LEGO reports seven scenes, without a verified matching
target roster. These recent rows are contextual, not same-protocol reproductions.
R averages all three seeds in each task. O uses our declared reference/source
contract. New NVOS evidence variants are reported as controlled extensions.}
\label{tab:nvos}
\end{table}

\paragraph{Multimodal UQIS.}
UQIS is the local ScanNet-UQIS9 v0.3.1 protocol; no external published row in
the cited papers evaluates these same 268 queries. Table~\ref{tab:uqis} therefore
reports all modalities and a retained regional-readout baseline. The difference
between regional readout and v10 measures our complete-decoder capability;
the optional selector has its own separate comparison. The matched v10-to-replacement comparison isolates the
response decision. The full
shared policy improves image IoU by 1.088 points and AP by 1.268 points. One
query improves and none are harmed. Text and both actual-point outputs remain
unchanged; the point rows verify preservation rather than learned improvements.
All reported outputs use the same policy, without choosing a different winning
variant for text, images or categories.
\begin{table}[ht]
\centering\small\resizebox{\linewidth}{!}{\input{tables/comparison_uqis.tex}}
\caption{UQIS matched comparison. T/I/P2/P3 denote text, image, 2D point and 3D
point. Actual selected-set IoU and continuous-response AP are distinct metrics.}
\label{tab:uqis}
\end{table}

\subsection{Historical component controls within our method}
Tables~\ref{tab:core-components}--\ref{tab:extent-components} connect v10's
capability to actual interventions. They are historical experiments on the
stated cohorts, not newly invented leave-one-out scores of the final system.

\begin{table}[ht]
\centering\small
\resizebox{\linewidth}{!}{\input{tables/mechanism_evidence.tex}}
\caption{Our matched identity, response-use and image-extent controls.
First two rows use LERF selected projection and semantic IoU respectively;
last two use actual UQIS image instance IoU. Different controls are non-additive.}
\label{tab:core-components}
\end{table}
Query admission changes candidate rejection with geometry, features and the set
solver fixed. Response-use control holds evidence, map and independent set
fixed while changing how admitted observations enter the semantic output; its
19.79-point gain is not attributed to a newly invented averaging equation.
Image local correspondence changes matching with the same teachers and extent
head. Complete-carrier reuse changes complete membership from the same identity
condition. These controls support different links of the same response process.

\begin{table}[ht]
\centering\small
\resizebox{\linewidth}{!}{\input{tables/provider_evidence.tex}}
\caption{Additional in-house component controls. Dual-native fusion uses three
held-out ScanNet rooms; point rows use UQIS actual-set IoU. Point supervision
changes observation distribution and sample counts, not only a formula.}
\label{tab:extent-components}
\end{table}
The 2.35-point category gain includes a second native teacher's capability.
Real-observation membership controls change training-observation distribution
and counts under the same learner. They must not be attributed solely to an
evidence-weighting rule or reported as zero-supervision results.

Four later full-cohort controls preserve these capabilities while testing
partial/complement permissions, available-domain normalization, scope agreement
and response solvers. No candidate produces common gains across all scored
outputs. Squared reconstruction reduces source loss in all 22 source-assisted
jobs, but downstream quality is mixed. These negative controls support keeping
effective typed operators; they are detailed in the diagnostic supplement.

\IfFileExists{tables/nvos_coverage.tex}{\input{sections/nvos_coverage.tex}}{}
\IfFileExists{tables/nvos_source_readout.tex}{\input{sections/nvos_source_readout.tex}}{}
\IfFileExists{tables/nvos_local_credit.tex}{\input{sections/nvos_local_credit_summary.tex}}{}

\subsection{Optional complete-response replacement controls}

All rows in Table~\ref{tab:ablations} retain the same teachers, complete
candidates, supervision rooms and output definitions. The query-blind arm returns only the fitted structural prior (64 trees,
versus 128 in the full model); it therefore also changes capacity. Permutation
instead uses the fitted full model and
shuffles all teacher scores jointly within each view. It preserves their
multisets and correlations. Removing the null constraint retains the same
64+64-tree capacity and total fit weight, replacing permuted training examples
with duplicates of aligned examples. Each trained control uses the same
calibration procedure. The zero-margin control uses the full fitted model and
changes only the acceptance margin. These are response-level ablations.

\begin{table}[ht]
\centering\small\resizebox{\linewidth}{!}{\input{tables/mechanism_ablations.tex}}
\caption{Matched controls on all ten ScanNet rooms and 268 UQIS queries.
S19/15/10 are mesh mIoU. Repair/harm counts concern actual UQIS set IoU;
unchanged actual points are included in the denominator. This is one coupled
policy comparison, not a collection of per-benchmark best runs.}
\label{tab:ablations}
\end{table}

The image gain over query-blind and permuted controls is respectively 0.725 and
1.088 points; descriptive nine-room bootstrap intervals are [0,2.174] and
[0,3.264]. They do not establish a robust query-specific advantage. The full
and no-null arms make identical choices in all 144 text/image/menu decisions,
so the null constraint has no demonstrated accuracy contribution. Removing the
positive margin improves text and image IoU, but reduces ScanNet mIoU, including
a 0.869-point loss on the ten-class menu. Retention prevents some damage while
also rejecting useful instance replacements. Neither the full policy nor the
zero-margin variant is promoted to replace the reference system.

\subsection{Candidate capacity and failure attribution}
For a fixed candidate pool, let $G_0$ be gain under zero-margin selection and
$G_\tau$ gain under calibrated retention. The exact diagnostic is
\[
G_{\rm pool}=R_{\rm rank}+(G_0-G_\tau)+G_\tau.
\]
The middle term is signed: it is negative when retention prevents damage.
This identity prevents attributing all unused candidate capacity to the margin.
For text, 15.395 points of available IoU gain decompose into 9.615 ranking regret,
5.780 retention cost and zero actual gain. For images, 9.492 available points
decompose into 8.224 ranking regret, 0.180 retention cost and 1.088 actual gain.
Thus image ranking is the dominant measured bottleneck; lowering the margin
cannot recover most of the unused capacity. Category mean-menu capacity is
only 0.913 points, and its negative retention cost indicates prevented damage.
These label-oracle diagnostics are computed after prediction and are never
automatic method results.

\paragraph{Five-task coverage and remaining evidence.}
The implemented v10 method covers every task. Positive component controls cover
LERF identity/response use, ScanNet native priors and UQIS identity/extent;
NVOS requires evidence and geometry attribution beyond scorer equivalence.
The optional selector has separate missing LERF/NVOS transfer cells. These two
kinds of completion must not be conflated. Independent teacher configurations
and unseen scenes remain needed for broader claims; no feature-dimension or
compression sweep is required for this response-focused paper.

\section{Discussion and limitations}
The experiments support a distinction between identifying evidence, authorizing
its spatial effect, and recovering a complete task output. Identity admission
and complete-membership controls show that these capabilities affect accuracy;
reconstruction and response-decision controls show that better source agreement
alone is insufficient. Their interventions and cohorts differ, so their gains
cannot be added. The contribution is the evidence-to-response design and its
measured behavior, built with native correspondence and established fusion
operations.

The common contract has typed implementations. Current evidence does not show
one admission rule or selector improving every task and teacher configuration.
The L11 null constraint has no demonstrated accuracy effect. The benchmark
cohort has informed development, and new NVOS variants are development
experiments rather than independent final tests. More teacher configurations
and unseen scenes are needed for broad transfer claims.

Source access, teacher strength and carrier geometry differ from some published
baselines. Gaussian versus mesh score differences are not method gains. LERF
selection uses projected GT, UQIS is a local protocol, and original NVOS v10 differs from
LUDVIG in target-RGB access and GT boundaries. The image-conditioned extension
explicitly changes that access, but keeps source-trained geometry and produces
view-conditioned updates. Early static image refinement improved IoU while
reducing AP; completing the original signed reference response subsequently
improves both. The N12--N14 variants nevertheless increase residuals against
their original source teachers on all eight tasks. Target-view improvements therefore
do not establish a universally improved or query-global 3D field.

The controls localize three distinct limitations. First, memory transport of
an incomplete original reference mask does not restore missing query scope.
Completing that reference before transport recovers the missing orchid flowers.
Second, source annotation agreement does not establish either complete extent
or accurate target boundaries: unconditional extra prompting can enlarge an
already adequate reference response. Source feasibility and predicted quality
are thus measured separately, with fixed decisions for every task. Third,
native-ray fitting reduces transport loss but cannot remove false positives
already present in the image response; fern and trex still expose internal
background leakage and thin-structure loss. N15 tests cross-view consistency:
completed-reference source-only fusion recovers 92.7461 IoU, but equally
weighting these sources with target evidence lowers N14's IoU by 0.1428 points
while increasing AP by 0.0346. Completed-source Huber residual falls from
0.004626 to 0.000748, yet foreground recall also falls. This controlled result
separates source-teacher consistency from task accuracy. The fusion branch is
stopped without a post-hoc weight search; it does not establish a universal
performance ceiling. The unchanged v10 threshold oracle (90.1098) diagnoses only that
original score field, not an upper bound on new spatial responses.

N16 shows that one extra target decode is also insufficient as a uniform
repair rule. Its local control modestly improves thresholded IoU but lowers AP;
the tradeoff already appears in direct image responses, before 3D fitting.
Its mean-IoU gain over N14 disappears when the correlated horns queries are
omitted together. This descriptive sensitivity is not a new validation split.
The new maximum IoU is retained in the tables rather than hidden, but the
extra decoding stage is not required by the recommended method. These stopped
branches support an empirical stopping decision for the current refinement
recipes, not a claim that no stronger evidence or architecture can help.
The subsequent A24 numerical audit exactly replays all 24 image responses.
Evaluating the same probability interpolation with float64 and a stable
complement-based ranking changes mean AP by less than 0.002 points and changes
no hard masks. The local image response remains 2.2124 AP points below matched
N14. Thus the suspected final sigmoid/interpolation precision effect does not
explain the observed ranking loss. This audit does not recompute 3D results.

The N11 geometric and v10-response anchors were nearly tied, so a distinct
geometric-admission advantage was not established. N13's paired probability
residual loses IoU on all eight tasks against its matched complete replacement.
Its exact no-intervention property is therefore insufficient evidence for a
better response rule. These alternatives remain diagnostic controls, not
mandatory stages of the current reference-completion pipeline. New admission
rules were developed after inspecting N12; they are explicitly development
results, without scene-name routing or target-score selection at inference.
Independent scorer checks reduce
known implementation concerns but do not prove that every possible bias has
been excluded.

\paragraph{Reproducibility.}
Versioned recipes, source/input hashes, per-task scores, protocol audits and
component controls bind claims to code and outputs. Teachers, maps and geometry
are referenced in place. New dense predictions are bounded and transient;
source snapshots, small models and reports persist. Table generation does not
run inference. Incomplete historical source archives remain explicitly marked.

\section{Conclusion}
Query-conditioned evidence fusion connects native evidence to geometric
responses while preserving distinct semantic and complete-selection outputs.
Our v10 system supplies the five-task implementation and main measurements.
Matched controls identify effective identity, response-use and extent operations;
negative controls show that better source fit alone is insufficient. NVOS
evidence experiments and optional response replacement are specific extensions,
with their effects separated from established method capability.

\input{sections/references.tex}
\end{document}
